% !TeX root = ../papers/nc-p3.tex
\section{負交換表示の全境界と二次曲面グラフへの特殊化}
\label{sec:negative-pairing-boundary}
前節と同じ $zD(z^2-1)\ne0$ を課し、$[a:b:c:d]$ は八例外直線を避ける。
\begin{theorem}\label{thm:negative-pairing-boundary}
この負交換表示の全正則点で、指定始点の同型の必要十分条件は
定理\ref{thm:negative-pairing-unmarked-quotient}の $H$ の軌道の一致である。
$q,t$ のうちちょうど一方が零なら、最小周期は $r=-p$ または $r=-zp$ で1、それ以外で2である。
両方が零なら、$(a+b)(a+zb)=0$ でちょうど二次曲面グラフとなり周期1、他は四直線のグラフを台に持ち周期2である。
この二次曲面条件も $r=-p$ または $r=-zp$ と同値である。
\end{theorem}
\begin{proof}[一方だけが零の場合の枠]
$B$ で $q,t$ を入れ替えられるので $d=a$, $t\ne0$ とする。前節の二次式 $P,Q$ を用い、
$\alpha=(a+b)(a+c)$, $\beta=(a+bz)(az+c)$ と置く。
$D(X_0-X_1)=(2a+b+c)X_2-(2az+bz^2+c)X_3$ と $\alpha X_2=\beta X_3$ が
$P=Q=0$ と同値である。$D\ne0$, $t\ne0$ により $\alpha,\beta$ は同時には消えない。
$X_0=X_1$ とすれば $Q=0$ が $X_2=zX_3$、$P=0$ が $(z-1)t X_3=0$ を与える。
従って $\alpha\beta\ne0$ なら二つの平面二次曲線が二点
$e_0\pm e_1$ で会い、それ以外なら平面二次曲線の二重化である。後者の被約二次曲線も滑らかである。

二次式はこの境界でも実際の小行列式を消す。具体的に、各小行列式は $P,Q$ の二次式係数結合である。
上位ビット別に混ざる項は前節の表示を使い、平方座標だけの五式は $P$ で $X_0$ を消去して $Q$ で割る。
$d=a$ では五つの商は全て $X_2,X_3$ の線型結合で分母は $D$ の冪だけである。
$Q\ne0$ と $Q=0,P\ne0$ の前節の消去は $q\ne0$ を使わず、
今度は $x_2=x_3=0$, $X_0=X_1$ の二点は曲線上に来る。孤立点は $x_0=x_1=0$, $X_2=zX_3$ の二点だけである。

二重二次曲線の場合にも $e_0\pm e_1$ は固有である。曲線の台で $x_2$ または $x_3$ が非零な点では、
零でない方の位置を使う前節の混合小行列式により、局所的にまず $Q$、続いて $P$ が生成される。
例えば $x_3=0,x_2\ne0$ の点では $(a+b,a+c)$ も $(c+d,b+d)$ も同時零を除いてあり、
$x_0,x_1$ の一方は非零である。この点において $P,Q$ の係数行列の階数が2だから
近傍のイデアルは $(P,Q)$ である。$x_2=0$ の場合は $z$ 付きの四係数で同様である。
二点 $e_0\pm e_1$ では $x_0x_1$ が単元で $P$ が生成される。
$P$ を法として全小行列式は $Q$ の極大イデアル倍であり、平方座標だけの五式は上の除算から分かる。
一方、$P$ による局所商は正則、$Q$ の像は非零である。
従って点スキームから純一次元部分を取る写像の核がこの両点で非零となる。
純部分は $(P,Q)$ そのものである。これは完全交叉で埋込み点がなく、既に二点以外で一致するためである。
よってこの二点は非Cohen--Macaulayな点として復元できる。

二つの二次曲線なら二つの節点と、二つの階数3の鉛筆要素の頂点を取る。
二重二次曲線なら上の二点を取り、二つの孤立点の割線と被約二次曲線の平面の交点を第三点、
その点の孤立二点に関する調和共役を第四点とする。いずれも得られる二組は
$\{e_0+e_1,e_0-e_1\}$ と $\{e_2,e_3\}$ であり、二組の役割は異なる。
奇位置では両組は $\{e_2+e_3,e_2-e_3\}$ と $\{e_0,e_1\}$ である。
実際、$BJ$ は $q=0$ に戻す表示で、位置の上位ビットを交換する。二節点側の元の座標直線では
指数の符号変更が二点を交換する。四位置で示した候補を六式に代入するといずれも評価核を持つので、
この直線と曲線の二交点が上の二点に等しい。孤立二点の座標も
\path{computations/verify_negative_boundary_vertices.py} で元の四辺から検算する。
\end{proof}

\begin{proof}[一方だけが零の場合の商]
復元した点を列とする基底でテンソルを書く。最大台58語、うち42語は
$z$ と非零定数だけで決まる。各位置で二組の各々を交換する $4^4$ 候補のうち、
この42語を最大台へ送るものは16通りだけで、どれも42語自体を保つ。
42語の整数行格子は階数11で、残りの左上係数 $2a,2b,2c$ の任意の一語を加えると、
全58語と同じ階数12の原始格子となる。三語の各対の指数差も42語の格子に入る。
よって $z$ と $[a:b:c]$ が順序付き枠から読め、射影対角安定群は一次元トーラスである。
より詳しくは四組の文字指数の核で42語の係数単項式から $z$ を得て、
三語の各比への補正にも42語の単元だけを用いる。従って射影的な三係数の比較は零成分へも延びる。

16通りのうち8通りは目標で $z'b'=c'$ を要求し、この層には戻らない。
この検査で $b'/c'$ を使った場合、分母が零になる点では分子も同じ元の係数に対応し、
両方とも零である。残りは $U,S$ の生成する四通りの $[a:b:c]$ の変換で、各々が二度現れる。
四通りは前節の実際の同型で、同じパラメータ作用を持つ二置換の差は、
各位置で節点対を交換する元の対角四元群の文字作用で実現する。
\path{computations/verify_negative_pairing_boundary.py} が全候補を照合する。
従って一般の射影安定群の連結成分数は2である。$z$ を保つパラメータ作用は
$[a:b:c]\mapsto[a:c/z:zb]$ だけなので、$a=0,r=0$ なら成分数4、他では2である。
$BJ$ は $[p:r:t]\mapsto[(z+1)p+2r:-2zp-(z+1)r:(z-1)t]$。
固定した $z$ で使える同型は $t$ の符号変更だけなので、最初の二成分の行列の
二固有線 $r=-p,-zp$ がちょうど周期1の条件となる。$(p,r)=(0,0)$ も両式を満たす。
\end{proof}

\begin{proof}[両方が零の場合]
$d=a$, $c=zb$ とする。この面では $D\ne0$, $z\ne0,\pm1$ のもとで八例外直線のどれにも乗らない。
$(a+b)(a+zb)\ne0$ なら $P,Q$ は $X_0-X_1,X_2-zX_3$ を張り、
前節の曲線外の点の計算から射影点スキームの台は四直線の四角形になる。
四つの節点は線型独立で、向かい合う組を入れ替えることができる。
$w^2=z$, $m=a/(wb)\in\mathbb P^1$, $X=(w+1)^2/(w-1)^2$, $Y=(m+1)/(m-1)$ と置く。
$X\ne0,\pm1,\infty$、$Y\ne0,\infty$ である。$b=0$ では $Y=1$ とする。
節点枠は偶位置で $e_0\pm e_1,we_2\pm e_3$、位置1で $we_0\pm e_1,e_2\pm e_3$、
位置3で $e_0\pm we_1,e_2\pm e_3$ である。これは上と同様に元の四辺への評価で確認できる。
この枠の28個の係数は全て非零で、文字倍率行列のSmith因子は $1^{11}$ である。
$b=0$ でも厚い左上を $b$ 倍、右下を $b^{-1}$ 倍して斉次座標に戻すと非零性は保たれる。
トーラス不変な単項式から $X-1,X,Y$ を読み取れる。
四位置で四角形の二面体置換を全て試すと、28語の台に戻るのは8通りで、
$(X,Y)$ には次の四通りが各二回現れる。
\[
 (X,Y),\quad(1/X,1/Y),\quad(X,-1/(XY)),\quad(1/X,-XY).
\]
これらが十分であることも確認する。$X,Y$ が等しい二表示では $w'=w$ または $1/w$ であり、
同じ $m$ である。後者は $U$ の作用で得られる。第二変換は平方根 $w$ の符号変更で、
第三変換は $B$ の作用である。同じ作用を持つ二置換は全位置の両対の入替えで異なり、
元の文字の下位ビットの符号により実現できる。$J$ は $(X,-Y)$ を与えるから
周期1の可能性は $XY^2=1$、すなわち先に除いた $(a+b)(a+zb)=0$ に限る。
この四角形の射影安定群は二次元トーラスを単位成分に持ち、成分数は $XY^2=-1$ で4、その他で2である。

残る二線は $B$ によって交換される。$a=1,b=-1,c=-z,d=1$ を取ると
$Q=x_0^2-x_1^2-x_2^2+zx_3^2$ で、第二位置を
$(x_0,x_1,x_2,x_3)\mapsto(-x_2,x_3,x_0,-x_1)$ により第一位置と同定できる。
六式の対称部分は $Q$ を張り、交代部分にある五平面の零化二ベクトルは
$\Omega=e_0\wedge e_1+e_2\wedge e_3$ である。非退化性により辺はスキームとして
滑らかな $Q$ の対角グラフである。各辺の基底を $B,J$ で戻すと同じ結論が得られる
（$J$ は $m\mapsto1/m$ で二線を交換する）。定理\ref{thm:smooth-quadric-graphs}の二平面の
面積形式の固有値比の平方は $z$ または $1/z$ である。実際、$Q$ の対角行列と
$\Omega$ の交代行列の積の特性多項式は $(h^2-1)(h^2-z)$ である。
従って正確な無印座標は $z+1/z$、周期は1である。ここでも同じ座標を持つ全点は
$B,U$ で同一視できる。この族で滑らかな二次曲面枝の座標 $-2$ 以外を全て実現するため、
その一変数枝の全体が負の四次元成分の閉包に含まれる。
以上の各形は射影点スキームの台の次元と孤立点数 $4,2,0$ で区別でき、層をまたぐ同型も尽きた。
\end{proof}

\begin{corollary}\label{cor:negative-exchange-opposite-ranks}
負交換表示の全域で反対階数は $16,15,12$ のいずれかである。階数15は商上の二因子で、
いずれも最小周期2であり、一位置移動で交換される。これらの接する階数12の全点は周期1で、
その指定始点の商は
\[
 (\mathbb A^1\setminus\{-2,2\})\times(\mathbb P^1\setminus\{2\})
\]
となる。二つの周期1平面は階数12に沿って交わる。
\end{corollary}
\begin{proof}
色0の反対ブロックの列を各二組の和・差に取ると、二行正方ブロック二つの直和になる。
一方は列 $(a,-d)^t,(1,1)^t$、他方は $(1,z)^t,((a+b+c+d)/D,-(az+bz^2+c+dz)/D)^t$ を持つ
（列の定数倍を除く）。色1でも二組の行の和・差で同じ分解があり、
最初の二列 $b,-c$ と $1,z$ の組および右下ブロックから同様に計算できる。
どの二行ブロックも零にならない。色2,3は両方とも行列式 $-4(z+1)^2$ である。
従って正確な階数は
\[
16-[p=0]-[r=0]-[2zp+(z+1)r=0]-[(z+1)p+2r=0].
\]
固定 $z$ で四平面は相異なり、共通の直線 $p=r=0$ を持つ。$B$ の作用は
$p=0$ と $2zp+(z+1)r=0$、$r=0$ と $(z+1)p+2r=0$ をそれぞれ交換し、$J$ は二組を交換する。
$z\ne0,\pm1$ では周期1平面がこれらに重ならない。

$p=r=0$ なら $d=-a,c=-zb$ で、八例外直線は $D\ne0$ の範囲でこの軌跡に会わない。
$h=t^2/(zq^2)\in\mathbb P^1$ とすれば $D\ne0$ は $h\ne1$ と同値である。
符号を除いた商で $B$ は $h\mapsto1/h$、$U$ は $h$ を保ち $z\mapsto1/z$ とする。
平方根の選択も符号変更で尽くされるから、商座標は $z+1/z$ と $h+1/h$ である。
$h=0,\infty$ は一方だけの対角差が消える境界から加わり、第二座標の無限遠の一点を与える。
\end{proof}
